\section{Related Work}

\subsection{Peptide--MHC Binding and Presentation Prediction}

Peptide--MHC binding prediction is a foundational problem in computational immunology. Established methods such as NetMHCpan \cite{reynisson2020netmhcpan}, MHCflurry \cite{odonnell2020mhcflurry}, and ACME \cite{hu2019acme} estimate binding affinity, binding rank, or peptide--MHC compatibility from peptide sequences and allele information. Pan-allelic predictors improve coverage by representing MHC molecules through pseudo-sequences that summarize the peptide-binding groove or through related allele features. This representation allows information sharing across well-characterized and data-scarce alleles \cite{nielsen2007netmhcpan,reynisson2020netmhcpan}. These models are essential for filtering candidate epitopes, but binding is only one stage of antigen presentation. A peptide with favorable binding characteristics may still be absent from the cell-surface immunopeptidome because its source protein is not expressed or degraded, the peptide is not generated or transported efficiently, or it is not detected experimentally. Binding scores therefore do not directly encode presentation-evidence preferences among tissues that share the same MHC restriction.

\subsection{Predicting Cell-surface Antigen Presentation (Antigen Presentation Prediction)}

Eluted-ligand data enable models to learn a broader presentation signal than binding measurements alone \cite{reynisson2020netmhcpan,odonnell2020mhcflurry}. General presentation predictors combine peptide sequence, MHC context, and, in some cases, antigen-processing or source-protein features to distinguish observed ligands from background peptides. Systems such as BigMHC \cite{albert2023bigmhc} illustrate how large immunopeptidomic collections and modern sequence models can improve ligand presentation prediction and candidate ranking. Their usual target, however, is whether a peptide is likely to be presented for an allele, sample, or pooled presentation setting. This target differs from the relative question studied here: given a particular tissue--MHC context, which of two peptides from the same parent protein has stronger tissue-conditioned presentation evidence? Consequently, a high general-presentation score need not imply a tissue-specific preference, and direct comparisons require a benchmark whose supervision reflects that distinction.

\subsection{Tissue- and Sample-Context Immunopeptidomics}

Multi-tissue immunopeptidomic atlases have revealed that ligand repertoires contain both broadly shared peptides and tissue-associated components \cite{marcu2021atlas,schuster2018mouse}. MHCflurry 2.0 incorporates antigen-processing features but does not use tissue identity \cite{odonnell2020mhcflurry}. More recent systems perform multi-allelic deconvolution, as in ImmuneApp \cite{xu2024immuneapp}, or combine peptide, MHC, peptide-flanking sequence, and optional gene-expression features, as in HLApollo \cite{thrift2024hlapollo}. These approaches demonstrate that presentation is context dependent, but they generally do not define a tissue--MHC combination as the supervised prediction unit or match every recorded positive and pseudo-negative on both MHC restriction and parent protein. Our formulation is therefore complementary to expression- and sample-aware prediction: it isolates a ranking problem conditioned on a known biological context that can later be combined with molecular measurements.

\subsection{Sharing Information across Tissue--MHC Combinations}

Some tissue--MHC combinations contain many observations, whereas others contain relatively few. Learning a completely separate model for every combination would waste information shared across tissues and alleles. Conversely, pooling all combinations into one model could blur allele-specific binding motifs and tissue-associated differences. Joint learning provides a middle ground: related combinations share a peptide representation while retaining separate prediction functions \cite{caruana1997multitask,ruder2017overview}. More structured methods allow different groups to use shared and specialized model components \cite{ma2018mmoe,tang2020ple}. TissuePMHC adopts a biologically structured sharing scheme: one model view learns patterns across all human combinations, while a second shares information only among tissues with the same HLA allele.

\subsection{Relation to This Work}

The principal distinction from previous work is the prediction target. General predictors estimate peptide binding to or presentation by an \plainMHC{}, whereas TissuePMHC ranks peptides matched by MHC restriction and parent protein within a specified tissue--MHC context. The human and mouse datasets define this comparison identically, although each species uses a separately trained model. We distinguish ordinary internal evaluation, which permits a peptide to recur in another fitting combination, from peptide-separated evaluation, in which a held-out peptide is absent from every fitting combination. MHCflurry and NetMHCpan are therefore used as tissue-independent peptide--MHC controls rather than as tissue-aware competitors.

Table~\ref{tab:closest-work} makes these distinctions explicit. Negative construction varies across implementations and datasets; the entries therefore summarize the published prediction design rather than asserting identical data provenance.

\begin{table*}[t]
\centering
\caption{Comparison with representative methods for peptide binding or presentation. ``Context'' denotes information used by the published method beyond the peptide and MHC.}
\label{tab:closest-work}
\fontsize{7.5}{8.5}\selectfont
\setlength{\tabcolsep}{2pt}
\begin{tabular}{@{}>{\raggedright\arraybackslash}p{1.8cm}>{\raggedright\arraybackslash}p{2.5cm}>{\raggedright\arraybackslash}p{2.3cm}>{\raggedright\arraybackslash}p{2.6cm}@{}}
\toprule
Method & Prediction target & Context & Negative/evaluation emphasis \\
\midrule
NetMHCpan 4.1 \cite{reynisson2020netmhcpan} & Binding and presentation learned from eluted ligands & MHC pseudo-sequence; no tissue identity & Pan-allelic prediction \\
\midrule
MHCflurry 2.0 \cite{odonnell2020mhcflurry} & Binding and presentation & Antigen-processing features; no tissue identity & Prediction across many alleles (pan-allelic general prediction) \\
\midrule
BigMHC \cite{albert2023bigmhc} & Presentation and immunogenicity & Peptide and MHC & Large-scale general prediction \\
\midrule
ImmuneApp \cite{xu2024immuneapp} & Epitope prediction and HLA-I deconvolution & Mono- and multi-allelic ligand data & General or sample-level evaluation \\
\midrule
HLApollo \cite{thrift2024hlapollo} & Peptide presentation & MHC, peptide flanks, and optional gene expression & Alternative negative-set definitions; pan-allelic evaluation \\
\midrule
TissuePMHC benchmark & Tissue--MHC presentation preference & Tissue and MHC identity & Peptide pairs matched by MHC and parent protein; ordinary and peptide-separated internal evaluation \\
\bottomrule
\end{tabular}
\end{table*}

Table~\ref{tab:closest-work} compares prediction designs rather than numerical performance.
